% Удиртгал

\chapter{Удиртгал}
\label{ChapterIntro}

Монголын загварын салбар сүүлийн жилүүдэд дизайнер, брэнд, коллекц, тайзны лүүк, редакцийн нийтлэл, зурагт материал, соёлын тайлбар зэрэг олон төрлийн контентоор баяжиж байна. Гэвч эдгээр мэдээлэл ихэвчлэн сошиал медиа, брэндийн тусдаа веб хуудас, богино мэдээ, зураг хадгалах сан, хувь хүний архив зэрэг олон эх сурвалжид тархсан байдаг. Ийм нөхцөлд уншигч, дизайнер, судлаач, редактор, админ хэрэглэгчид нэг дизайнерын коллекц, нийтлэл, зураг, материал, холбоотой брэндийг нэгдсэн байдлаар хайж үзэхэд хүндрэлтэй.

Энэхүү дипломын ажил нь Монголын загварын мэдээллийг сэтгүүлийн редакцийн бүтэцтэйгээр нэгтгэн харуулах, нийтлэл болон коллекцийн мэдээллийг удирдах, зурагт контентыг архивлах, хэрэглэгч хайлт хийх, хадгалах, мөн архивын өгөгдөлд тулгуурласан AI туслахаас лавлах боломжтой Anoce/Mongolian Fashion Archive веб платформыг зохиомжлон хөгжүүлэхэд чиглэнэ. Уг платформ нь зөвхөн танилцуулгын сайт бус, харин public fashion magazine, дижитал архив, admin CMS, content repository, глобал хайлт, хэрэглэгчийн хадгалах/follow үйлдэл, RAG chatbot зэрэг олон давхаргатай мэдээллийн платформ юм.

\section*{Зорилго}
Энэхүү дипломын ажлын зорилго нь Монголын загварын дизайнер, коллекц, лүүк зураг, редакцийн нийтлэл, ангилал, зурагт медиа файл болон хэрэглэгчийн үйлдлийг нэгтгэн удирдах, хайх, үзүүлэх, AI туслахаар лавлах боломжтой загварын сэтгүүлийн веб платформ зохион бүтээж, хэрэгжүүлэхэд оршино.

\section*{Зорилт}
Энэхүү зорилгыг хэрэгжүүлэхийн тулд дараах зорилтуудыг дэвшүүлэв:
\begin{itemize}
    \item Загварын сэтгүүл, онлайн редакцийн платформ, CMS болон дижитал архивын онолын суурь ойлголтыг судлах;
    \item Дотоодын болон гадаадын ижил төрлийн загварын, редакцийн, архивын, визуал хайлтын платформуудыг харьцуулан шинжлэх;
    \item Уншигч, редактор, админ хэрэглэгчийн шаардлага, use case, бизнес процесс, өгөгдлийн логик бүтцийг тодорхойлох;
    \item Дизайнер, коллекц, лүүк, нийтлэл, ангилал, медиа файл, хэрэглэгч, эрх, хадгалсан контент, RAG document зэрэг өгөгдлийн нэгжийг загварчлах;
    \item Next.js, React, TypeScript, Tailwind CSS, CouchDB, Supabase ашиглан public фронтэнд, admin CMS, content API, authentication, хайлт болон AI chatbot-ийн үндсэн хэрэгжилтийг хийх;
    \item Платформын архитектур, API зохиомж, өгөгдлийн сангийн зохиомж, UI/UX зохиомж, аюулгүй байдлын зохиомжийг диаграм, хүснэгтээр баримтжуулах;
    \item Хөгжүүлсэн платформын үндсэн функц, фронтэнд/бэкенд ажиллагаа, authentication, responsive UI, хайлт болон AI туслахын туршилтыг хийж үр дүнг үнэлэх.
\end{itemize}

\section*{Платформ хөгжүүлэх үндэслэл}
Монголын загварын контент нь зураг, дизайнерын profile, коллекц, лүүк, нийтлэл, чиг хандлагын тайлбар, материалын мета өгөгдөл зэрэг олон хэлбэртэй тул ердийн static site эсвэл нийтлэг мэдээний сайт дангаараа хангалтгүй. Загварын сэтгүүлийн платформд visual-first layout, дизайнер-коллекц-лүүк хоорондын холбоо, редакцийн нийтлэл, зурагт медиа файл, ангилал, хайлт, хэрэглэгчийн оролцоо, редакцийн workflow зэрэг онцлог шаардлагууд зэрэгцэн ордог.

Платформ хөгжүүлэх үндэслэлийг дараах байдлаар тодорхойлов:
\begin{enumerate}
    \item \textbf{Контентыг нэгтгэх хэрэгцээ:} Монголын загварын мэдээлэл олон эх сурвалжид тархсан тул уншигч нэг дороос нийтлэл, зураг, коллекц, дизайнерын мэдээлэл үзэх боломж сул байна.
    \item \textbf{Редакцийн удирдлагын хэрэгцээ:} нийтлэл, дизайнер, коллекц, медиа файл байнга нэмэгдэх тул фронтэнд код өөрчлөхгүйгээр админ болон редактор контент удирдах CMS шаардлагатай.
    \item \textbf{Зурагт архивын хэрэгцээ:} загварын сэтгүүлийн үнэ цэнэ нь зөвхөн текст мэдээнд бус, лүүк зураг, материал, улирал, брэндийн түүх, визуал навигацид оршино.
    \item \textbf{Хайлтын болон AI лавлагааны хэрэгцээ:} архив томрох тусам хэрэглэгч дизайнер, коллекц, материал, чиг хандлага, нийтлэлийн түлхүүр үгээр хайж, асуултаа энгийн хэлээр асуух шаардлага нэмэгдэнэ.
\end{enumerate}

\section*{Шинэлэг тал}
Энэхүү платформын шинэлэг тал нь Монголын загварын сэтгүүлийн контентыг зөвхөн мэдээний жагсаалт байдлаар бус, редакцийн CMS, дизайнер/коллекцийн архив, глобал хайлт, хэрэглэгчийн хадгалах/follow үйлдэл, AI туслахтай нээлт бүхий нэг платформ болгон зохион байгуулсанд оршино. Үүнд:
\begin{itemize}
    \item Дизайнер, коллекц, лүүк, нийтлэл, ангилал, медиа файлын өгөгдлийг нэг content model-д холбосон;
    \item Public archive, editorial, designer directory болон collection detail хуудсуудыг зураг төвтэй UX-ээр зохион байгуулсан;
    \item CouchDB content repository ба Supabase authentication/role/bookmark/RAG layer-ийг хослуулсан;
    \item Уншигч, редактор, админ гэсэн эрхийн түвшинд суурилсан CMS бүтэц хэрэгжүүлсэн;
    \item Live archive болон curated RAG dataset ашиглан Монголын загварын контентын талаар хариулах Anoce chatbot-ийг нэгтгэсэн.
\end{itemize}

\section*{Ач холбогдол}
Платформ хэрэгжсэнээр уншигч Монголын дизайнеруудын коллекц, лүүк, материал, нийтлэл, чиг хандлагын тайлбарыг нэг газраас үзэх боломжтой болно. Редактор, админ хэрэглэгчид нийтлэл, дизайнер, коллекц, зурагт медиа файл нэмэх, засах, хянах, нийтлэх ажиллагаагаа нэг платформд төвлөрүүлнэ. Судалгааны хувьд уг ажил нь Монголын загварын салбарын мэдээллийг веб технологи, өгөгдлийн сан, UX дизайн, CMS workflow, AI туслахтай уялдуулсан дижитал сэтгүүл-архивын жишээ болно.

\section*{Хамрах хүрээ}
Энэхүү дипломын ажлын хүрээнд public home, archive, designers, editorial, journal, careers, press, article detail, collection detail, search overlay, profile/bookmark, бүртгэл устгах, admin dashboard, article CRUD, collection CRUD, designer CRUD, asset upload, review/publish workflow, user role gate, admin-only invite/role update, content API болон AI chatbot-ийн үндсэн хэрэгжилтийг авч үзнэ. Production payment gateway, бүрэн хэмжээний analytics, олон хэлний localization, бүх Монголын загварын албан ёсны бүрэн каталогийг нэг дор бүрдүүлэх зорилго энэ ажлын хүрээнд хамаарахгүй.
